%% sections/02-related-work.tex

\section{Related Work}
\label{sec:related}

\subsection{Narrative Structure and Campaign Design}

Classical narrative structure has been applied to game design through frameworks
ranging from Freytag's pyramid~\cite{campbell1949hero} to three-act structure as
adapted for interactive media~\cite{ryan2017narrative}. These frameworks provide
a gross structural shape (introduction, escalation, climax) without specifying what
design mechanisms ensure the escalation and climax are reached in multi-session play.
The campaign spine's 4-hint delivery schedule is a more granular structural
specification: it plants and delivers content at designed intervals, creating a
cadenced convergence trajectory rather than leaving progression to improvisation.

The closest structural antecedent in practitioner literature is the \textit{mystery
plot structure}: clues planted early, delivered progressively, and combined at the
finale to enable a solution~\cite{mystery2012}. The campaign spine adapts this to
emotional and thematic content rather than informational content. The hints are not
clues about ``who did it'' but pieces of insight about what the campaign's central
question means and how it might be answered. This distinction is fundamental: a
mystery's resolution is an intellectual event; the spine's finale is an
emotional and volitional event.

Published TRPG campaign design resources offer structural guidance that is
largely advisory rather than validated. Adventure path design (as practiced in
Pathfinder and similar products) specifies adventure sequences with designed
connective tissue between modules but stops short of specifying explicit
arc-completion conditions~\cite{paizo2021}. Apocalypse World's \textit{fronts}
system~\cite{baker2010} provides a threat-escalation architecture that creates
structural pressure toward confrontation, but not a designed completion state.
Burning Wheel's beliefs/instincts/traits (BIT) architecture~\cite{crane2002}
gives each character an explicit belief that the campaign should answer, but
the campaign's arc is not structurally specified beyond the individual
character level. The campaign spine formalizes what these systems gesture toward:
an explicit, cross-character arc architecture with designed completion conditions.

\subsection{Quest Design and Campaign-Level Structure}

Quest design in digital games has been studied from the perspectives of goal
structure~\cite{dormans2010}, player motivation~\cite{tychsen2006role}, and
procedural generation~\cite{shaker2016procedural}. This body of work operates
primarily at the single-quest level --- the equivalent of an adventure, not a
campaign. The campaign spine addresses a level above quest design: it governs
the multi-quest arc rather than individual quests. Adventure design within the
campaign is quest-level work; the spine is campaign-level architecture.

The gap between these levels has received some attention in digital game narrative
research. \citet{riedl2016} describe AI planning approaches for narrative sequencing
that generate quest chains with goal coherence. However, goal coherence
(``each quest contributes to a goal'') differs from emotional arc completion
(``the party's engagement with the campaign's central question produces a
satisfying answer by the finale''). The spine's emphasis on earned completion ---
where the party's specific actions during specific sessions determine which
ending they reach --- is not addressed by goal-directed narrative planning.

\citet{fan2019} examine strategies for structuring long-form narrative generation,
finding that narrative coherence across extended sequences requires mechanisms
for tracking and referencing established facts. The campaign spine's
campaign-permanent fact system (events from played sessions that constrain all
future sessions) directly implements this requirement at the design level.

\subsection{Story Completion in Interactive Drama and AI Narrative Systems}

\citet{mateas2003faade} address story completion in interactive drama through
\textit{drama management}: an AI director monitors the dramatic curve and steers
player behavior toward designed outcomes. The campaign spine takes a structurally
different approach. Rather than steering behavior during play, it designs explicit
conditions under which the campaign can end at any of several quality levels.
Players who engage deeply with the spine's central question achieve Route D or E;
players who engage less deeply achieve Routes A--C. The campaign always completes;
the quality of completion varies with engagement level. This is completion with
variance, not completion through direction.

This is closer to \citet{ryan2017narrative}'s principle that narratives should be
completable by any player, with the depth of completion varying with engagement
level. The spine operationalizes this principle through its route structure: every
campaign has a minimum-engagement completion path (Route A or C) that requires
no special party behavior, and progressively harder paths (Routes D and E) that
require sustained engagement with the spine's design intent.

AI narrative generation research has produced systems capable of generating
story events~\cite{ammanabrolu2021, fan2019}, but these systems do not address
the campaign-level problem of sustained emotional arc across multiple sessions
with persistent state and player-driven outcomes. The Marathon workshop addresses
this problem in an AI-simulated context: an LLM serves as DM across all 7 sessions,
with campaign-permanent facts logged to persistent files that constrain all
future sessions. The spine is the design artifact that makes this possible at
scale.

\subsection{Trust, Token Systems, and Progress Mechanics}

Token and reputation systems in digital games (Mass Effect's Paragon/Renegade
scale, Dragon Age's companion approval ratings) provide mechanics for tracking
player relationship with the narrative's ethical and emotional axes. These systems
share the spine's trust-accumulation logic but differ in implementation: digital
game reputation systems operate continuously (every action increments or decrements
a counter), while the spine's token system is event-driven (specific named
party actions earn specific tokens). The event-driven design is deliberate: it
ensures that token-earning events are recognizable moments in the narrative
rather than undifferentiated incremental accumulation.

Tabletop precedents include Dungeon World's bond system~\cite{koebel2012} and
Powered by the Apocalypse moves, which make specific player actions trigger
specific consequences. The spine's token-earning conditions are closer to moves
than to continuous reputation tracking: each token has an explicit condition,
and meeting that condition is the meaningful act, not the token itself.

The spine's route-threshold design (different routes available at different token
counts) is related to branching narrative structures in digital games~\cite{ryan2017narrative},
but differs in that routes are not pre-authored paths through different content.
All routes use the same adventure content; what differs is which ending scene
the party reaches and what the PC-authored deliverable contains. This is
\textit{depth-of-engagement branching}, not content-path branching.

\subsection{Convergent Finale Structure}

The simultaneous delivery of all four hints in the finale session --- a structural
pattern observed in all 7 Marathon campaigns --- has antecedents in classical
dramatic theory. Aristotle's \textit{recognition and reversal} (\textit{anagnorisis}
and \textit{peripeteia}) as simultaneous events describes the structural moment
when understanding and consequence arrive together~\cite{aristotle1996}. The spine's
simultaneous hint delivery is this structure operationalized: hints planted
separately across the campaign converge at the finale moment, producing a
recognition event for the party that enables the optimal ending.

Screenwriting theory describes the equivalent as the ``obligatory scene'' ---
the scene the audience has been waiting for since the premise was established ---
and emphasizes that its emotional impact depends on all prior scenes having done
their work~\cite{field2005}. The spine's hint schedule is the mechanism for
ensuring the prior work is done: each hint is a piece of the finale that the
party carries forward from its planting session.
